\chapter{Caracterización Cayley–Pontryagin de los ceros críticos y reconstrucción efectiva del campo espectral de Weil}
\label{chap:hmt83-doble-circulo-campo-espectral}

Una vez localizada la recta crítica por la factorización positiva, la
transformada de Cayley publica las alturas como fases de un primer círculo. El
entrelazador de Pontryagin las transporta después al círculo nonádico sin
identificar retorno de fase con retorno de estado. El teorema del doble círculo
y la cuantización de las órbitas se articulan con la reconstrucción efectiva
del campo espectral. La exposición conserva explícitamente las hipótesis, los
residuos de Galerkin, el cociclo de memoria, el transporte torcido y el
entrelazador que identifica ambas realizaciones sin confundir el índice de la
vuelta con la altura espectral.

\Needspace{9\baselineskip}
\section{Anulación de la función completada y cierre espectral sobre la recta crítica}
La imagen convencional de los ceros de Riemann suele comenzar demasiado tarde. Se presenta la función zeta, se prolonga analíticamente y se pregunta en qué puntos se anula. El lector ve una función compleja y una colección de puntos misteriosos distribuidos sobre una recta.\par
La lectura HMT cambia el orden de la pregunta. No comienza mirando dónde se hace cero la publicación final. Comienza reconstruyendo el estado que esa publicación ha contraído: los relojes de las potencias primas, el canal gamma, el canal polar, las dos hojas, su orientación, el acarreo, la memoria, la frontera, el terminal y la conexión nonádica que los transporta conjuntamente.\par
Sólo después aparece la función completada y se pregunta por sus ceros.\par
La diferencia es radical: un cero deja de ser una ausencia y pasa a ser un acontecimiento de cierre.\par
\Needspace{7\baselineskip}
\subsection{Anulación escalar y persistencia del estado enriquecido}
Si\par
\[
\xi(\rho)=0,
\]
el cero es verdadero. No es aproximado, ficticio ni decorativo. La función completada se anula exactamente en \(\rho\).\par
Pero lo que se anula es una publicación escalar. No desaparece el punto \(\rho\), no desaparece su multiplicidad, no desaparece la estructura prima que participa en la fórmula explícita y no desaparece el estado enriquecido que precede a la publicación.\par
Ésta es una distinción elemental y, sin embargo, decisiva:\par
\[
\text{valor publicado igual a cero}
\quad\not\Rightarrow\quad
\text{estado generador inexistente}.
\]
Una representación útil es una franja oscura producida por interferencia. La oscuridad no significa que ninguna onda haya llegado. Significa exactamente lo contrario: han llegado contribuciones cuya relación de fase es tan precisa que su suma visible se cancela.\par
Los ceros no triviales pueden leerse de esta manera: como nodos de interferencia aritmética global.\par
No son nodos de una sola secuencia de primos ni pertenecen individualmente a un primo. En la publicación completa intervienen simultáneamente:\par
\begin{itemize}
\item los relojes de las potencias primas,
\end{itemize}
\[
\ell_{p,k}=k\log p;
\]
\begin{itemize}
\item el canal arquimediano gamma;
\item el canal polar, que conserva la normalización y retira los polos y ceros triviales;
\item las dos hojas especulares;
\item la frontera que la proyección visible no muestra.
\end{itemize}
Un cero es una cancelación exacta de esa publicación conjunta. Su existencia revela una coordinación, no un vacío.\par
Por tanto, los ceros constituyen \textbf{nodos espectrales de cancelación aritmética}: la anulación de \(\xi\) expresa una coordinación global exacta entre las contribuciones de la publicación conjunta.\par
\Needspace{7\baselineskip}
\subsection{Fijación de la recta crítica por la involución funcional \(s\mapsto1-\bar s\)}
La función completada conserva las simetrías\par
\[
\rho\longmapsto1-\rho,
\qquad
\rho\longmapsto\overline\rho.
\]
Combinándolas aparece la involución\par
\[
\rho\longmapsto\rho^\sharp
:=
1-\overline\rho.
\]
Si\par
\[
\rho=\beta+i\gamma,
\]
entonces\par
\[
\rho^\sharp
=
1-\beta+i\gamma.
\]
La reflexión conserva la altura \(\gamma\), pero intercambia \(\beta\) y \(1-\beta\). Su conjunto de puntos fijos es:\par
\[
\rho=\rho^\sharp
\quad\Longleftrightarrow\quad
\beta=1-\beta
\quad\Longleftrightarrow\quad
\beta=\frac12.
\]
Por tanto,\par
\[
\operatorname{Fix}(\sharp)
=
\left\{
\frac12+i\gamma:\gamma\in\mathbb R
\right\}.
\]
La recta crítica no es solamente “la mitad” por una preferencia geométrica. Es la costura donde las dos hojas funcionales dejan de aparecer como posiciones distintas y se convierten en una sola posición autodual.\par
Fuera de la recta, un cero aparece ligado a una órbita de hasta cuatro puntos:\par
\[
\rho,\qquad
1-\rho,\qquad
\overline\rho,\qquad
1-\overline\rho.
\]
Sobre la recta crítica, la reflexión funcional conjugada ya no envía el cero a otra posición lateral:\par
\[
1-\overline\rho=\rho.
\]
El cero cierra consigo mismo.\par
Éste es el primer sentido preciso del círculo de cierre: no se trata únicamente de que dos mitades sumen uno, sino de que la órbita especular encuentra un punto fijo.\par
Sin embargo, esto todavía no demuestra la hipótesis de Riemann. La igualdad\par
\[
\frac12=1-\frac12
\]
identifica el eje; no obliga por sí sola a que todos los ceros estén sobre él.\par
La obligación procede de la positividad.\par
\Needspace{7\baselineskip}
\subsection{Cuadraticidad del producto cruzado bajo clausura de la involución funcional}
La forma espectral completa puede escribirse como\par
\[
\mathscr I_{\mathrm{GW}}(f,g)
=
\sum_{\rho\in\mathscr Z}
\Phi_f(\rho)\,
\overline{\Phi_g(1-\overline\rho)},
\]
donde \(\mathscr Z\) es el multiconjunto de ceros no triviales, conservando sus multiplicidades.\par
Esta fórmula determina cómo se emparejan los ceros: la evaluación situada en \(\rho\) se combina con la evaluación situada en su reflejo \(\rho^\sharp=1-\overline\rho\).\par
Si el cero está fuera de la recta, la expresión empareja dos puntos diferentes:\par
\[
\Phi_f(\rho)\,
\overline{\Phi_f(\rho^\sharp)}.
\]
Eso es un producto cruzado. No es necesariamente un módulo cuadrado y, sobre una clase de pruebas suficientemente rica, puede producir direcciones de signo negativo.\par
Si el cero está en la recta crítica, entonces\par
\[
\rho^\sharp=\rho
\]
y el mismo término se convierte en\par
\[
|\Phi_f(\rho)|^2.
\]
La forma completa pasa a ser:\par
\[
\mathscr I_{\mathrm{GW}}(f,f)
=
\sum_{\rho\in\mathscr Z}
|\Phi_f(\rho)|^2.
\]
Ahora el espejo no empareja dos posiciones desplazadas. Cada punto se empareja consigo mismo y la publicación se convierte en una suma de cuadrados.\par
Ésta es la lectura conceptual de la demostración: HMT construye la forma de Weil aguas arriba como una energía de frontera,\par
\[
\mathscr I_{\mathrm{GW}}
=
E_R^*E_R
=
\mathscr L_R^*\mathscr L_R
\succeq0.
\]
No se declara positiva porque los ceros estén sobre la recta. Se construye como cuadrado antes de usar los ceros.\par
Después, el criterio reversible de Weil dice que esa positividad sobre toda la clase separadora de funciones de prueba es equivalente a que todos los ceros no triviales satisfagan\par
\[
\operatorname{Re}\rho=\frac12.
\]
Por tanto, la positividad no “atrae” dinámicamente a los ceros hacia el centro. Los ceros no se desplazan desde una posición inicial. Lo que ocurre es más fuerte: cualquier configuración fuera del eje fijo produciría una dirección negativa detectable, mientras que el complemento HMT es materialmente una norma cuadrada y no puede adquirirla.\par
La recta crítica es el único lugar compatible con ambas publicaciones.\par
\Needspace{7\baselineskip}
\subsection{Caracterización espectral de los ceros}
Bajo esta lectura, un cero no trivial puede describirse simultáneamente de cuatro maneras compatibles.\par
Es un \textbf{cero analítico verdadero}, porque\par
\[
\xi(\rho)=0.
\]
Es un \textbf{punto fijo del espejo}, porque\par
\[
\rho=1-\overline\rho.
\]
Es un \textbf{átomo espectral positivo}, porque en la forma de Weil aparece como\par
\[
|\Phi_f(\rho)|^2.
\]
Y es un \textbf{nodo global de los relojes primos y del borde arquimediano}, porque no pertenece a un canal aislado: surge de la publicación coordinada de los términos primo, gamma y polar.\par
El cero es nulo como amplitud y positivo como localización espectral. No hay contradicción. Son dos lectores diferentes.\par
La función vale cero en \(\rho\), pero \(\rho\) continúa apareciendo, con toda su multiplicidad, dentro de la medida espectral positiva. El cero no ha sido borrado por ser cero. Al contrario: se convierte en uno de los lugares que sostienen la suma positiva.\par
Éste es probablemente el cambio interpretativo más importante:\par
\begin{quote}
El cero no es el lugar donde la estructura desaparece. Es el lugar donde la estructura completa alcanza una cancelación visible y deja una marca espectral.\par
\end{quote}
\Needspace{7\baselineskip}
\subsection{Identificación de la altura de un \(0\) no trivial con una frecuencia espectral real}
La normalización empleada en el criterio de Weil es\par
\[
\Phi_f(s)
=
\widehat f\!\left(
i\left(s-\frac12\right)
\right).
\]
Si\par
\[
\rho=\frac12+i\gamma,
\]
entonces\par
\[
i\left(\rho-\frac12\right)
=
-\gamma
\]
y, por tanto,\par
\[
\Phi_f(\rho)
=
\widehat f(-\gamma).
\]
La altura \(\gamma\) se convierte literalmente en una frecuencia real de Fourier.\par
Así, la recta crítica puede leerse como una recta de frecuencias: la coordenada real \(\tfrac12\) fija el equilibrio especular y la coordenada \(\gamma\) registra la frecuencia del nodo.\par
Esto explica por qué el lenguaje espectral es apropiado. Los ceros no son solamente coordenadas complejas; son frecuencias en las que la publicación aritmética completa manifiesta su condición nodal.\par
Estas alturas admiten la interpretación de las frecuencias nodales de un sistema global: puntos donde una amplitud observable se cancela mientras la estructura que produce esa cancelación permanece activa.\par
Pero hay que proteger esta frase contra una exageración. De aquí no se sigue automáticamente que cada \(\gamma\) sea un autovalor energético de un Hamiltoniano físico concreto. Para eso habría que construir ese Hamiltoniano, su dominio y la correspondencia espectral. Lo que está demostrado es la lectura Fourier–Mellin, la forma positiva completa y la localización de los ceros.\par
Se trata, con rigor, de frecuencias espectrales. No se convierten, sin otro entrelazador, en niveles físicos de energía.\par
\Needspace{7\baselineskip}
\subsection{Parametrización circular conforme de la recta crítica}
Tu intuición del círculo no es sólo metafórica. Existe una transformación de Möbius exacta:\par
\[
\tau(s)
=
\frac{s-1}{s}.
\]
Se cumple:\par
\[
\operatorname{Re}s=\frac12
\quad\Longleftrightarrow\quad
|\tau(s)|=1.
\]
Por tanto, la recta crítica compactificada por su punto del infinito se transforma en la circunferencia unidad.\par
Si\par
\[
s=\frac12+it,
\]
entonces\par
\[
\tau(s)\in S^1.
\]
Al escribir\par
\[
\tau=e^{i\theta},
\]
la inversa es\par
\[
s=\frac1{1-\tau}
=
\frac12+\frac{i}{2}\cot\frac{\theta}{2}.
\]
La altura \(t\) de la recta se convierte en una posición angular sobre el círculo. Cuando \(t\) recorre toda la recta, la carta recorre la circunferencia compactificada. Los dos extremos \(t\to-\infty\) y \(t\to+\infty\) se encuentran en el punto \(\tau=1\), aunque lo hacen desde orientaciones opuestas.\par
También el espejo adquiere una forma circular:\par
\[
\tau(1-\overline s)
=
\frac1{\overline{\tau(s)}}.
\]
La involución espectral se convierte en reflexión respecto de la circunferencia unidad, cuyo conjunto fijo es precisamente el círculo.\par
Así que la expresión “círculo de cierre” posee un contenido exacto:\par
\begin{itemize}
\item la recta crítica compactificada es un círculo;
\item el espejo funcional se convierte en reflexión respecto de ese círculo;
\item los ceros no triviales son puntos marcados sobre esa frontera;
\item cada punto está fijado por la reflexión que antes intercambiaba las dos hojas.
\end{itemize}
El círculo no genera los ceros ni determina sus alturas. Toda la recta se transforma en el círculo, mientras que la ecuación\par
\[
\xi(\rho)=0
\]
selecciona sus puntos espectrales.\par
Tampoco significa que los ceros estén igualmente espaciados o que haya nueve por vuelta. El círculo es la geometría compactificada del locus crítico; la distribución de los ceros sobre él conserva su propia aritmética.\par
\Needspace{7\baselineskip}
\subsection{Círculo nonádico y holonomía de fase}
Hay un segundo círculo, anterior al analítico: la base cíclica de las nueve fases.\par
La conexión nonádica no consta de nueve filtros ni de nueve canales independientes. Es una única holonomía:\par
\[
\Gamma_9
=
\operatorname{Hol}_{C_9}
(\nabla_{\mathrm{TPK}}).
\]
Su ley de retorno es\par
\[
g(K+9)=g(K),
\]
pero simultáneamente\par
\[
q_{\mathrm{mem}}
(\Gamma_9\widehat x)
=
q_{\mathrm{mem}}(\widehat x)+1
\]
y, en general,\par
\[
\Gamma_9\widehat x
\neq
\widehat x.
\]
La fase vuelve; el estado no se reinicia.\par
Visto desde la base, se ha recorrido un círculo y se ha regresado a la misma fase. Visto desde el estado enriquecido, se ha acumulado una vuelta de memoria. La imagen más fiel para el lector es:\par
\begin{itemize}
\item un círculo en la proyección;
\item una hélice en el estado completo.
\end{itemize}
Desde arriba, la hélice parece cerrar. Desde dentro, cada vuelta ocupa una altura distinta.\par
La parte literal no es la imagen de la hélice, sino las identidades matemáticas: retorno de fase, incremento de memoria y cambio general del estado.\par
Esta estructura fue comprobada en 396 niveles de refinamiento por canal. El certificado registra 43 vueltas nonádicas completas y 387 comparaciones \(K,K+9\) por canal: la fase coincide y el estado cambia en todas ellas. No es una manera poética de hablar de periodicidad; es una separación ejecutada entre retorno visible y continuidad genealógica.\par
\Needspace{7\baselineskip}
\subsection{El espejo nonádico y el espejo de Riemann}
Tu intuición es correcta: en ambos casos aparece un espejo. Pero conviene decir con exactitud cómo se relacionan.\par
El espejo interno nonádico intercambia las dos componentes laterales, asociadas a las lecturas \(\pi/e\), mientras conserva el eje \(\varphi\) y el agregado correspondiente. La involución de Riemann es\par
\[
s\longmapsto1-\overline s.
\]
No son literalmente el mismo mapa, porque tienen dominios y codominios distintos. Igualarlos sin el entrelazador sería confundir la arquitectura interna con su publicación analítica.\par
La afirmación correcta es más interesante: el entrelazador espectral–polar publica la arquitectura two-way del TPK como la involución funcional de la función completada. El espejo nonádico es anterior; el espejo de Riemann es su realización en el codominio analítico.\par
En ambos casos existe:\par
\begin{itemize}
\item una pareja de hojas;
\item un eje que permanece fijo;
\item una transformación que intercambia las posiciones laterales;
\item una costura donde las dos publicaciones coinciden;
\item una memoria que no desaparece al efectuar la identificación.
\end{itemize}
La recta \(\Re s=\tfrac12\) es la costura visible de ese espejo en el plano complejo.\par
\Needspace{7\baselineskip}
\subsection{Cierre residual y fase nonádica}
La expresión “los nueves son falsos ceros” contiene una intuición muy precisa si distinguimos el entero, la clase residual y el estado completo.\par
Para \(n>0\), la carta positiva utiliza\par
\[
\rho_9(n)
=
1+((n-1)\bmod9).
\]
Equivalentemente,\par
\[
\rho_9(n)
=
\begin{cases}
n\bmod9,
&
n\not\equiv0\pmod9,\\[1mm]
9,
&
n\equiv0\pmod9.
\end{cases}
\]
Por tanto, cuando un número pertenece a la clase nula módulo nueve, la carta positiva no publica \(0\): publica \(9\).\par
Al mismo tiempo se conserva\par
\[
\kappa_9(n)
=
\frac{n-\rho_9(n)}9
\]
y la reconstrucción\par
\[
n
=
9\kappa_9(n)
+
\rho_9(n).
\]
El nueve es, así, el representante positivo de la clase residual cero. Es cero en el cociente, pero no es el entero cero y no es un estado vacío.\par
El nueve dice: “la vuelta se ha completado”. El cociente dice cuántas coronas o vueltas se han acumulado. Si se sustituye todo por el residuo \(0\), se conserva la congruencia, pero se pierde la diferencia entre ausencia, cierre y memoria.\par
Ése es el sentido riguroso de “falso cero”:\par
\begin{quote}
El nueve es un cero para la sombra residual, pero sería falso interpretarlo como aniquilación del estado. Es un cero publicado con memoria de cierre.\par
\end{quote}
Por ejemplo, \(27\), \(72\) y \(729\) poseen sombra residual cero y raíz digital positiva nueve. Sin embargo, no son el mismo número ni tienen la misma genealogía. Uno puede proceder de una reversión, otro de una evaluación diferente y otro de una potencia. El residuo no reconstruye la ruta.\par
Esto conecta profundamente con los ceros de Riemann, pero no los identifica.\par
Un cero de \(\xi\) es un cero analítico auténtico. No es el número nueve. No existe aquí una afirmación de que cada cero no trivial corresponda individualmente a una fase nueve concreta.\par
La relación estructural es ésta:\par
\[
\text{una publicación puede anularse}
\quad\text{sin que se anule su genealogía}.
\]
El nueve es una falsa ausencia aritmética. El cero de zeta es una falsa ausencia ontológica. Ambos son ceros verdaderos dentro de su lector, pero ninguno significa que no haya estructura detrás.\par
\Needspace{7\baselineskip}
\subsection{Separación tipada de los ceros de \(\zeta\), \(\xi\) y el determinante espectral}
Para que la fuerza de la narración no borre los tipos matemáticos, conviene separar tres ceros.\par
\Needspace{5\baselineskip}
\subsubsection{\(0\) analítico como anulación de la función en el dominio espectral}
\[
\xi(\rho)=0.
\]
Es un punto espectral de la función completada. Conserva posición, altura, multiplicidad y simetrías.\par
\Needspace{5\baselineskip}
\subsubsection{\(0\) terminal como objeto universal del morfismo de cierre}
En la telescopía HMT aparece\par
\[
(A^NE_R)^*(A^NE_R)
=
q^{2N}
\mathcal B_{\mathrm{sp},R}^{\mathrm{str}}
\longrightarrow0,
\qquad
q=\frac59.
\]
Es la extinción de un resto después de que toda su contribución haya sido transferida a una suma de cuadrados positivos.\par
\Needspace{5\baselineskip}
\subsubsection{Representante positivo de la clase residual nula: \(n\equiv0\pmod 9\) y \(\rho_9(n)=9\)}
\[
n\equiv0\pmod9,
\qquad
\rho_9(n)=9.
\]
Es la clase nula modular publicada como representante positivo, acompañada por su cociente y memoria.\par
No son el mismo objeto. Pero comparten una gramática:\par
\[
\text{anulación visible}
+
\text{conservación estructural}.
\]
\Needspace{7\baselineskip}
\subsection{Agotamiento del objeto terminal bajo la filtración cofinal}
La demostración conserva en cada profundidad una identidad finita:\par
\[
\mathcal B_{\mathrm{sp},R}^{\mathrm{str}}
=
\sum_{n=0}^{N-1}
(DA^nE_R)^*(DA^nE_R)
+
(A^NE_R)^*(A^NE_R).
\]
El último término es el terminal. No se elimina para obtener una suma positiva más bonita. Permanece presente en cada etapa finita.\par
Sólo después se prueba:\par
\[
(A^NE_R)^*(A^NE_R)
=
q^{2N}
\mathcal B_{\mathrm{sp},R}^{\mathrm{str}}
\longrightarrow0.
\]
Esto permite otra lectura del cero: el terminal llega a cero porque su contenido ha sido expresado y contabilizado nivel tras nivel. Cada parte perdida por la contracción visible ha quedado registrada en una coordenada de memoria positiva.\par
El resto se extingue, pero no ha sido borrado.\par
Es como vaciar un depósito trasladando exactamente su contenido a una secuencia de recipientes numerados. Al final, el depósito inicial está vacío, pero nada ha desaparecido: existe un ledger completo de la transferencia.\par
Aquí el cero significa \textbf{transferencia terminada}.\par
\Needspace{7\baselineskip}
\subsection{Información genealógica del continuo codificada por los ceros críticos}
La recta crítica\par
\[
\frac12+i\mathbb R
\]
es una copia desplazada de la recta real. Desde HMT, esa recta no es el objeto inicial: es una publicación arquimediana del continuo genealógico.\par
Los ceros forman un conjunto discreto dentro de esa publicación continua. No perforan la recta ni interrumpen el continuo. Constituyen un divisor espectral: puntos aislados, cada uno con multiplicidad, sobre una sombra continua.\par
Esta convivencia de continuo y discreto es esencial. La recta ofrece el codominio continuo; la genealogía determina las marcas discretas que aparecen sobre ella.\par
La imagen que propondría al lector es la de una membrana continua atravesada por nodos espectrales. Los nodos no rompen la membrana. Hacen visible su estructura interna.\par
En términos holográficos, el cero es una marca de frontera producida por una organización global del interior. No pertenece a un único primo, igual que un píxel de una imagen holográfica no pertenece a un único punto del objeto representado. Cada cero responde al conjunto completo de relaciones primo–gamma–polar.\par
Esto no significa que cada cero contenga por sí solo toda la aritmética de manera recuperable. Significa que su posición no puede atribuirse a una entrada local aislada: es una respuesta global del sistema.\par
\Needspace{7\baselineskip}
\subsection{Correspondencia entre cuantización espectral, frecuencias reales y realización física del operador}
Los ceros se interpretan como nodos de una transferencia aritmética global:\par
\begin{itemize}
\item la coordenada \(\gamma\) actúa como frecuencia;
\item la recta crítica es la superficie de equilibrio especular;
\item la forma de Weil es la energía positiva del complemento;
\item la conexión nonádica conserva la memoria de cada vuelta;
\item el terminal mide lo que todavía no ha sido expresado;
\item el cero visible señala una cancelación perfecta de la amplitud publicada.
\end{itemize}
No se identifican con partículas, agujeros, estados de vacío o masas, ni con niveles energéticos materiales sin construir antes un operador físico adicional.\par
La interpretación física válida es estructural: nodo, frecuencia, espejo, transferencia, memoria y cierre. La ontología física concreta requeriría su propio transductor.\par
\Needspace{7\baselineskip}
\subsection{Corolarios de localización, multiplicidad y simetría de la caracterización espectral}
Antes de la resolución, la pregunta parecía ser:\par
\begin{quote}
¿Por qué unos puntos definidos por una función compleja parecen alinearse misteriosamente sobre una recta?\par
\end{quote}
Después de la resolución, la pregunta puede formularse de otro modo:\par
\begin{quote}
¿Cómo publica la estructura espectral–polar una energía de frontera que sólo puede permanecer positiva cuando cada cero se identifica con su reflejo?\par
\end{quote}
La respuesta es:\par
\begin{itemize}
\item las columnas positiva y negativa proceden del mismo estado;
\item el estado de frontera \(E_R\) conserva lo que la proyección no muestra;
\item la diferencia de los dos Grams es
\end{itemize}
\[
E_R^*E_R;
\]
\begin{itemize}
\item la telescopía conserva el terminal hasta demostrar su extinción;
\item la forma de Weil completa coincide con ese cuadrado;
\item el criterio reversible obliga a que cada cero sea un punto fijo del espejo.
\end{itemize}
La recta ya no es una coincidencia visual. Es el conjunto fijo de la involución compatible con la positividad construida.\par
\Needspace{7\baselineskip}
\subsection{Hipótesis y codominio exactos del teorema de localización espectral}
Esta lectura permite afirmar:\par
\[
\{\text{ceros no triviales}\}
\subset
\left\{
\frac12+i\gamma:\gamma\in\mathbb R
\right\}.
\]
Y, mediante la carta de Möbius,\par
\[
\left\{
\frac12+i\gamma
\right\}
\longrightarrow
S^1.
\]
Por tanto, todos los ceros no triviales son puntos marcados del círculo crítico.\par
No se ha construido aquí una biyección individual:\par
\[
\{
\text{ceros no triviales}
\}
\longleftrightarrow
\{
\text{vueltas nonádicas concretas}
\}.
\]
No se afirma que el primer cero sea la primera vuelta, que haya nueve ceros por ciclo o que sus alturas se generen repitiendo una fase. La resolución demuestra localización, positividad, cierre especular y pertenencia al círculo. Una cuantización individual de las alturas sería un desarrollo adicional y necesitaría su propio mapa.\par
Tampoco deben confundirse el espejo interno \(\pi/e\), la involución funcional de Riemann y una eventual simetría física. La relación entre ellos es de publicación tipada, no de identidad nominal.\par
\Needspace{7\baselineskip}
\subsection{Teorema de cierre espectral mediante la factorización positiva de Guinand–Weil}
La interpretación operatoria se resume mediante las propiedades siguientes:\par
\begin{quote}
Los ceros no triviales de Riemann son anulaciones exactas de la amplitud analítica que conservan la genealogía de su construcción. Cada cero determina una frecuencia global, un marcador espectral y un punto autodual de la involución funcional. La compactificación de la recta crítica produce el círculo fijo común de las dos hojas. La conexión nonádica distingue el retorno de fase del reinicio del estado: el cierre circular conserva el avance de memoria. La clase nonádica \(9\equiv0\pmod 9\) representa el cierre visible acompañado por la información interna del estado enriquecido.\par
\end{quote}
La anulación espectral constituye así una forma concentrada de cierre operatorio y conserva el objeto genealógico subyacente.\par
Esta interpretación físico-matemática acompaña a los resultados demostrados sin sustituirlos y no afirma una correspondencia individual entre ceros y vueltas nonádicas mientras no se construya el entrelazador correspondiente.\par

\section{Forma de Weil, traslación y cuantización de las alturas espectrales}
El objeto correcto no es una vuelta desnuda, sino una órbita helicoidal enriquecida:\par
\[
\boxed{
\text{cero no trivial}
\longleftrightarrow
\text{átomo espectral ponderado}
\longleftrightarrow
\text{carácter helicoidal nonádico}.
}
\]
Esta formalización actúa después del continuo completo y no redefine ni separa sus cinco características.\par
\Needspace{7\baselineskip}
\subsection{Construcción del espacio de Hilbert mediante la forma de Weil}
Se toma\par
\[
\mathcal D=C_c^\infty(\mathbb R),
\qquad
\langle f,g\rangle_W
:=
\mathscr I_{\mathrm{GW}}(f,g).
\]
En la resolución HMT, esta forma no nace de la lista de ceros. Se construye primero desde los canales primo–potencia, gamma y polar, y queda factorizada por la pérdida nonádica:\par
\[
\mathscr I_{\mathrm{GW}}
=
E^*E
=
\mathscr L^*\mathscr L
\succeq0.
\]
Se define\par
\[
\mathcal N_W
=
\{f\in\mathcal D:
\mathscr I_{\mathrm{GW}}(f,f)=0\},
\]
y completamos:\par
\[
\mathcal H_W
=
\overline{\mathcal D/\mathcal N_W}.
\]
Éste es el espacio de Weil positivo construido desde el lado aritmético. Los ceros todavía no han sido utilizados para definirlo.\par
Al mismo tiempo, la telescopía proporciona el espacio nonádico de pérdida\par
\[
\mathcal K_{\mathrm{loss}}
=
\overline{\operatorname{Ran}\mathscr L},
\]
y la igualdad de normas define una unidad canónica\par
\[
W_{\mathrm{loss}}:
\mathcal H_W
\longrightarrow
\mathcal K_{\mathrm{loss}},
\qquad
W_{\mathrm{loss}}[f]=\mathscr Lf.
\]
No estamos colocando después una copia artificial de la forma de Weil dentro del TPK: el propio espacio de pérdida de la prueba es una realización de su Hilbert positivo.\par
\Needspace{7\baselineskip}
\subsection{Operador de traslación independiente de la localización previa de los ceros}
En el dominio global \(\mathcal D\), no en una ventana fija, se considera\par
\[
(T_tf)(x)=f(x-t).
\]
La invariancia puede comprobarse antes de diagonalizar la fórmula en los ceros. Con la convención\par
\[
\widehat{T_tf}(z)
=
e^{-izt}\widehat f(z),
\]
se tiene\par
\[
h_{T_tf,T_tg}(z)=h_{f,g}(z),
\qquad
\check h_{T_tf,T_tg}(a)=\check h_{f,g}(a).
\]
Por tanto, todos los términos primo, gamma y polar permanecen invariantes:\par
\[
\mathscr I_{\mathrm{GW}}(T_tf,T_tg)
=
\mathscr I_{\mathrm{GW}}(f,g).
\]
La traslación conserva \(\mathcal N_W\) y desciende a un grupo unitario:\par
\[
U_t[f]=[T_tf].
\]
Después de la localización crítica ya demostrada, su continuidad fuerte sigue de\par
\[
\|U_t[f]-[f]\|_W^2
=
\sum_\gamma
m_\gamma
\left|e^{it\gamma}-1\right|^2
\left|\widehat f(-\gamma)\right|^2
\longrightarrow0.
\]
El teorema de Stone produce entonces un operador autoadjunto \(H_W\):\par
\[
U_t=e^{itH_W}.
\]
Este operador se ha definido mediante traslaciones y la forma primo–gamma–polar. No se ha definido escribiendo previamente una matriz diagonal con los ceros.\par
\Needspace{7\baselineskip}
\subsection{Identificación espectral de las alturas de los ceros no triviales}
Como la hipótesis de Riemann ya está demostrada, cada cero no trivial tiene la forma\par
\[
\rho_\gamma=\frac12+i\gamma.
\]
Si \(\Gamma_\xi\) es el conjunto de alturas distintas y \(m_\gamma\) su multiplicidad, se introduce la medida de conteo ponderada\par
\[
\mu_\xi
=
\sum_{\gamma\in\Gamma_\xi}
m_\gamma\,\delta_\gamma.
\]
La fórmula espectral de Weil se convierte en\par
\[
\mathscr I_{\mathrm{GW}}(f,g)
=
\sum_{\gamma\in\Gamma_\xi}
m_\gamma\,
\widehat f(-\gamma)
\overline{\widehat g(-\gamma)}.
\]
Por tanto,\par
\[
\mathcal A[f](\gamma)
=
\widehat f(-\gamma)
\]
define una isometría\par
\[
\mathcal A:
\mathcal H_W
\longrightarrow
L^2(\Gamma_\xi,\mu_\xi).
\]
Además, no queda una parte espectral escondida fuera de su rango. El rango de \(\mathcal A\) es cerrado y se conserva por todos los multiplicadores\par
\[
M_t a(\gamma)=e^{it\gamma}a(\gamma).
\]
La proyección sobre un átomo individual se obtiene mediante el promedio espectral\par
\[
P_\gamma
=
\operatorname*{s-lim}_{T\to\infty}
\frac1{2T}
\int_{-T}^{T}
e^{-it\gamma}M_t\,dt.
\]
La separación de evaluaciones de las transformadas de funciones de prueba garantiza que, para cada \(\gamma\), existe algún \(f\) con\par
\[
\widehat f(-\gamma)\neq0.
\]
Por tanto, \(P_\gamma\operatorname{Ran}\mathcal A\neq0\). Como el rango es reductor, contiene el eje atómico correspondiente. Todos los ejes atómicos pertenecen al rango y generan densamente \(L^2(\mu_\xi)\). En consecuencia,\par
\[
\boxed{
\mathcal H_W
\simeq
L^2(\Gamma_\xi,\mu_\xi)
}
\]
y\par
\[
\boxed{
\mathcal A H_W\mathcal A^{-1}
=
M_\gamma.
}
\]
Así pues,\par
\[
\boxed{
\sigma(H_W)=\Gamma_\xi,
}
\]
sin alturas adicionales, sin alturas omitidas y sin componente continua.\par
Esto es ya una cuantización espectral individual: las alturas de los ceros son exactamente el espectro puntual de un operador definido desde la forma aritmética positiva.\par
La multiplicidad exige una precisión. La forma escalar registra\par
\[
\mu_\xi(\{\gamma\})=m_\gamma.
\]
Es decir, conserva exactamente el peso algebraico del cero. No produce automáticamente \(m_\gamma\) vectores distinguibles: si un cero fuera múltiple, la evaluación en ese mismo punto se repite con peso \(m_\gamma\). La correspondencia canónica es entre el divisor de ceros y los átomos espectrales ponderados. Separar internamente las copias exigiría una forma positiva de jets o demostrar simplicidad.\par
\Needspace{7\baselineskip}
\subsection{Transformación conforme de la recta crítica en el círculo espectral con conservación de alturas}
Se usa la normalización\par
\[
\tau(s)=\frac{s-1}{s}.
\]
Entonces\par
\[
\Re s=\frac12
\quad\Longleftrightarrow\quad
|\tau(s)|=1.
\]
Para\par
\[
\rho_\gamma=\frac12+i\gamma
\]
se obtiene\par
\[
\tau_\gamma
=
\tau(\rho_\gamma)
=
\frac{-\frac12+i\gamma}{\frac12+i\gamma}
=
\frac{\gamma+\frac i2}{\gamma-\frac i2}
\in S^1.
\]
Ésta es exactamente la transformada de Cayley del operador \(H_W\):\par
\[
\mathcal C_W
=
\left(H_W+\frac i2I\right)
\left(H_W-\frac i2I\right)^{-1}.
\]
Sobre el átomo \(\gamma\),\par
\[
\mathcal C_WP_\gamma
=
\tau_\gamma P_\gamma.
\]
La correspondencia no pierde información. Su inversa es\par
\[
\rho_\gamma=\frac1{1-\tau_\gamma}.
\]
Si\par
\[
\tau_\gamma=e^{i\theta_\gamma},
\]
entonces\par
\[
\boxed{
\gamma
=
\frac12\cot\frac{\theta_\gamma}{2}.
}
\]
Por tanto, la altura completa está contenida en el ángulo. La compactificación de la recta conserva la localización de cada cero.\par
Además,\par
\[
\tau_{-\gamma}
=
\overline{\tau_\gamma}
=
\tau_\gamma^{-1}.
\]
Los ceros conjugados corresponden a orientaciones inversas del mismo círculo. Éste es el espejo convertido en una identidad angular exacta.\par
La correspondencia demostrada es\par
\[
\boxed{
\sum_\rho m_\rho[\rho]
\longleftrightarrow
\sum_\gamma m_\gamma
[\tau_\gamma,P_\gamma].
}
\]
No depende de ordenar previamente los ceros por altura.\par
\Needspace{7\baselineskip}
\subsection{Monodromía de cada \(0\) sobre el círculo espectral}
Transportamos ahora el operador circular al espacio nonádico:\par
\[
\mathcal C_{\mathrm{loss}}
=
W_{\mathrm{loss}}\,
\mathcal C_W\,
W_{\mathrm{loss}}^{-1}.
\]
En la telescopía HMT,\par
\[
q=\frac59
\]
es la contracción correspondiente a una vuelta completa. Se define la realización espectral de esa vuelta:\par
\[
\mathcal M_{\mathrm{spec}}
=
q\,\mathcal C_{\mathrm{loss}}.
\]
Si \(v_\gamma\) pertenece al átomo espectral \(P_\gamma\), entonces\par
\[
\boxed{
\mathcal M_{\mathrm{spec}}^n v_\gamma
=
q^n\tau_\gamma^n v_\gamma.
}
\]
Ésta es la órbita helicoidal del cero.\par
Su radio es\par
\[
\|\mathcal M_{\mathrm{spec}}^n v_\gamma\|
=
q^n\|v_\gamma\|,
\]
mientras su ángulo acumulado es\par
\[
n\theta_\gamma.
\]
La parte radial conserva la extinción terminal; la parte angular conserva la identidad individual del cero. Y ambas satisfacen la telescopía:\par
\[
\|v_\gamma\|^2
=
(1-q^2)
\sum_{k=0}^{N-1}
q^{2k}\|v_\gamma\|^2
+
q^{2N}\|v_\gamma\|^2.
\]
El terminal\par
\[
q^{2N}\|v_\gamma\|^2
\longrightarrow0
\]
no borra la altura: la altura permanece en el carácter unitario \(\tau_\gamma^n\).\par
Ésta es la lectura correcta de la vuelta:\par
\Needspace{8\baselineskip}
\begin{center}
\small
\begin{tabularx}{\textwidth}{@{}YY@{}}
\toprule
Dato & Qué registra \\
\midrule
\(n\) & cuántas vueltas completas han transcurrido \\
\(q=5/9\) & cuánto terminal queda después de cada vuelta \\
\(\tau_\gamma\) & qué cero individual transporta la órbita \\
\(m_\gamma\) & el peso algebraico de ese cero \\
memoria/carry & por qué volver a una fase visible no devuelve el mismo estado \\
\bottomrule
\end{tabularx}
\end{center}
Por eso no hay que asignar el primer cero a la primera vuelta. Cada cero atraviesa toda la escalera de vueltas con su propio carácter.\par
Una imagen pedagógica adecuada sería la de un disco que gira: el número de vueltas no identifica la nota que suena. La vuelta \(n\) dice cuánto tiempo ha transcurrido; la frecuencia identifica el modo. Aquí \(n\) es el contador, \(q\) es la atenuación y \(\tau_\gamma\) es la firma espectral.\par
\Needspace{7\baselineskip}
\subsection{Localización del campo espectral de Weil mediante el entrelazador nonádico de prolongación}
La construcción anterior vive realmente en\par
\[
\overline{\operatorname{Ran}\mathscr L},
\]
el espacio de pérdida producido por la telescopía nonádica. No es una analogía externa añadida a la prueba de Riemann.\par
Pero debe distinguirse de la holonomía combinatoria nativa. El TPK contiene también\par
\[
M_9=qV_9,
\qquad
V_9^*V_9=I.
\]
La identidad de defecto determina \(q\), pero no determina la parte unitaria \(V_9\). Un mismo balance\par
\[
M_9^*M_9=q^2I
\]
es compatible con muchas fases unitarias diferentes. La telescopía, por sí sola, no puede decir cuál de ellas contiene las alturas.\par
La identidad entre las órbitas espectrales recién construidas y las órbitas del ledger combinatorio APP–TRIT–TPK queda establecida por el entrelazador que se construye, sobre los estados cilíndricos enriquecidos, en el teorema Cayley--Pontryagin de este mismo capítulo:\par
\[
J:
\mathcal K_{\mathrm{loss}}
\longrightarrow
\mathcal K_{\mathrm{TPK}}^{\mathrm{unit}}
\]
tal que\par
\[
\boxed{
J\mathcal C_{\mathrm{loss}}
=
V_9J.
}
\]
Además, \(J\) debe conservar cilindro, orientación, carry, ruta, memoria, frontera y multiplicidad espectral.\par
La igualdad anterior no se postula ni se recibe de la realización clásica: resulta del cociclo entero de memoria, del transporte cilíndrico torcido por \(\mathcal C_{\mathrm{Weil}}\), de la proyectividad de los pesos y del encaje isométrico \(J_k\) construido más abajo. Al iterarla se obtiene\par
\[
(qV_9)^nJv_\gamma
=
q^n\tau_\gamma^nJv_\gamma,
\]
y entonces cada cero queda identificado con el carácter unitario de una órbita nativa del TPK. La identidad conserva cilindro, orientación, acarreo, ruta, memoria, frontera y multiplicidad espectral, de modo que la correspondencia no depende de una enumeración externa cero--vuelta.\par
\Needspace{7\baselineskip}
\subsection{Algoritmo espectral de cálculo independiente de la introducción previa de los ceros}
El resultado no se limita a representar una lista ya conocida. El operador \(H_W\) se ha definido desde traslaciones y desde la forma primo–gamma–polar. Su resolvente puede escribirse como\par
\[
(H_W-iI)^{-1}
=
i\int_0^\infty
e^{-t}U_{-t}\,dt.
\]
Por tanto, para funciones de prueba \(f_j,f_k\),\par
\[
\left\langle
(H_W-iI)^{-1}[f_j],[f_k]
\right\rangle_W
=
i\int_0^\infty
e^{-t}
\mathscr I_{\mathrm{GW}}(T_{-t}f_j,f_k)\,dt.
\]
El miembro derecho se calcula desde los canales primo, gamma y polar. No necesita alimentar una tabla de alturas.\par
Como el espectro es discreto y \(|\gamma|\to\infty\), el resolvente es compacto. Si \(P_N\) son proyecciones finitas que convergen fuertemente a la identidad,\par
\[
P_N(H_W-iI)^{-1}P_N
\longrightarrow
(H_W-iI)^{-1}
\]
en norma.\par
Sus autovalores convergen hacia\par
\[
\lambda_\gamma
=
\frac1{\gamma-i},
\]
y de ellos se recupera\par
\[
\gamma=i+\lambda_\gamma^{-1}.
\]
Esto abre una ruta computacional genuina: formar matrices finitas usando solamente la forma aritmética construida directamente desde los datos y recuperar las alturas como límites espectrales.\par
Las cotas efectivas de truncamiento y la elección de bases optimizadas pertenecen a la realización numérica del resolvente compacto; refinan su eficiencia computacional sin condicionar el teorema espectral ni el entrelazamiento operatorio ya construido. No hay una enumeración tautológica del tipo \(n=N(\gamma_n)\): los coeficientes de las aproximaciones proceden de la forma aritmética y sus autovalores convergen al espectro del operador.\par
\Needspace{7\baselineskip}
\subsection{Cuantización espectral y entrelazamiento nonádico}
La correspondencia exacta se formula en el siguiente nivel tipado:\par
\begin{quote}
La correspondencia probada no identifica ordinalmente el cero número \(j\) con la vuelta nonádica número \(j\), ni distribuye nueve ceros por ciclo, ni impone una repetición periódica de alturas. Identifica exactamente el divisor de ceros no triviales con los átomos ponderados del generador de traslaciones de la forma de Weil. La transformada de Cayley convierte cada altura en un carácter unitario único; la realización nonádica de pérdida lo transporta como una órbita helicoidal \(q^n\tau_\gamma^n\); y el transporte torcido \(\widetilde U_k\), junto con \(J_k\), entrelaza ese carácter espectral con la holonomía combinatoria nativa.\par
\end{quote}
La interpretación operatoria resultante es la siguiente:\par
\begin{quote}
Cada cero determina el carácter angular individual de una historia espectral cuyo iterado atraviesa sucesivas vueltas nonádicas. El índice de vuelta cuantifica la profundidad de memoria y el cero fija el carácter unitario que gobierna su evolución angular.\par
\end{quote}
En consecuencia, el resultado establece una vía operatoria efectiva: demuestra la cuantización espectral individual de las alturas, construye su órbita nonádica canónica y la identifica, mediante \(J_k\), con el carácter de memoria que actúa sobre todas las vueltas del ledger TPK. La identificación es armónica y operatoria; no depende de numerar los ceros con las vueltas.\par
\paragraph{Dominio y construcción del campo espectral de Weil}
La forma positiva \(\mathscr I_{\mathrm{GW}}=\mathscr L^*\mathscr L\), la contracción \(q=5/9\), la telescopía y la vuelta \(M_9=qV_9\) constituyen la estructura matemática de partida. Sobre ella se construyen el espacio de Hilbert de Weil, el grupo de traslaciones, su generador de Stone, la correspondencia de Cayley y las órbitas \(q^n\tau_\gamma^n\). La aproximación del resolvente directamente desde la forma aritmética define su realización computacional. El cociclo de memoria y el transporte torcido construidos a continuación proporcionan el entrelazador exacto con la holonomía combinatoria nativa.

\section{Teorema Cayley–Pontryagin del doble círculo}
No existe una correspondencia “cero número \(j\) = vuelta nonádica número \(j\)”. Lo que sí existe es algo más profundo: cada cero determina un carácter unitario de la memoria vertical del ledger, y la órbita completa de ese carácter es exactamente la órbita espectral del cero después de la transformación de Cayley.\par
\Needspace{7\baselineskip}
\subsection{Teorema del doble círculo}
La conexión nonádica posee una memoria vertical entera\par
\[
q_{\mathrm{mem}}:\operatorname{Obj}(\mathrm{TPK})\longrightarrow\mathbb Z.
\]
Para cualquier camino admisible \(\gamma:x\to y\), se define el cociclo\par
\[
m(\gamma)
=
q_{\mathrm{mem}}(y)-q_{\mathrm{mem}}(x).
\]
Es aditivo:\par
\[
m(\gamma_2\gamma_1)
=
m(\gamma_2)+m(\gamma_1),
\]
y para toda vuelta completa de la conexión,\par
\[
m(\Gamma_9)=1.
\]
Por iteración,\par
\[
m(\Gamma_9^n)=n.
\]
Éste es el dato decisivo. La fase visible vuelve después de nueve pasos, pero el estado no vuelve al mismo punto: la memoria ha avanzado una unidad.\par
Por otro lado, la demostración de Riemann construye primero la forma positiva de Weil y su espacio GNS:\par
\[
\mathcal H_W
=
\overline{\mathcal D/\ker \mathscr I_{\mathrm{GW}}}.
\]
La traslación construida desde la forma aritmética produce un generador autoadjunto \(H_W\). Sin introducir ceros como datos, se forma su transformada de Cayley:\par
\[
\mathcal C_{\mathrm{Weil}}
=
\left(H_W+\frac{i}{2}I\right)
\left(H_W-\frac{i}{2}I\right)^{-1}.
\]
Como \(H_W\) es autoadjunto,\par
\[
\mathcal C_{\mathrm{Weil}}^*
\mathcal C_{\mathrm{Weil}}=I.
\]
Por tanto, \(\mathcal C_{\mathrm{Weil}}\) es unitario y su espectro pertenece a \(S^1\).\par
Después del reconocimiento espectral, no antes, una altura\par
\[
\rho_\gamma=\frac12+i\gamma
\]
produce la fase\par
\[
\tau_\gamma
=
\frac{\gamma+i/2}{\gamma-i/2}
=
\frac{\rho_\gamma-1}{\rho_\gamma},
\qquad
|\tau_\gamma|=1.
\]
La inversión es exacta:\par
\[
\rho_\gamma=\frac{1}{1-\tau_\gamma},
\qquad
\gamma
=
\frac12\cot\left(\frac{\arg\tau_\gamma}{2}\right).
\]
Así aparecen dos realizaciones circulares obtenidas por lectores diferenciados del mismo estado enriquecido:\par
\[
\left\{\Re s=\frac12\right\}\cup\{\infty\}
\overset{\tau(s)=(s-1)/s}{\cong}
S^1,
\]
y\par
\[
\widehat{\mathbb Z_{\mathrm{mem}}}
\cong S^1.
\]
El primero es el círculo de Cayley de la recta crítica. El segundo es el dual de Pontryagin de la memoria vertical del ledger. El entrelazador demuestra que no son dos círculos meramente parecidos: son las dos realizaciones del mismo operador de fase.\par
\Needspace{7\baselineskip}
\subsection{Isometría de transporte torcido \(\widetilde U_k\) sobre estados cilíndricos, pesos proyectivos y memoria de Weil}
Sea \(Z_k\) el conjunto de estados cilíndricos enriquecidos en profundidad \(k\), con medida proyectiva de soporte completo \(\mu_k\), y sea\par
\[
\tau_{k+1,k}:Z_{k+1}\longrightarrow Z_k
\]
el truncamiento.\par
Para un descendiente \(z'\) de \(z\), se define\par
\[
w_k(z'|z)
=
\sqrt{\frac{\mu_{k+1}(z')}{\mu_k(z)}}.
\]
La proyectividad da\par
\[
\sum_{\tau_{k+1,k}z'=z}w_k(z'|z)^2=1.
\]
Sobre\par
\[
\mathscr K_k
=
\ell^2(Z_k)\otimes\mathcal H_W
\]
se define el transporte torcido\par
\[
\widetilde U_k(e_z\otimes v)
=
\sum_{\tau_{k+1,k}z'=z}
w_k(z'|z)\,
e_{z'}\otimes
\mathcal C_{\mathrm{Weil}}^{\,m(z\to z')}v.
\]
En forma matricial,\par
\[
[\widetilde U_k]_{z',z}
=
\mathbf 1_{\{\tau z'=z\}}
\sqrt{\frac{\mu_{k+1}(z')}{\mu_k(z)}}
\,
\mathcal C_{\mathrm{Weil}}^{\,m(z\to z')}.
\]
Esta fórmula conserva materialmente todos los descendientes. No suma caminos distintos como si fueran el mismo: ledger, orientación, carry, frontera, memoria, ruta y terminal continúan presentes en la etiqueta \(z'\).\par
La isometría se demuestra directamente. Los descendientes de padres diferentes son ortogonales; para cada padre, los pesos tienen suma cuadrática uno; y cada potencia de \(\mathcal C_{\mathrm{Weil}}\) es unitaria. Por tanto,\par
\[
\widetilde U_k^*\widetilde U_k=I.
\]
La composición también es exacta:\par
\[
\widetilde U_{\ell,j}\widetilde U_{j,k}
=
\widetilde U_{\ell,k},
\]
porque los pesos telescopan y el cociclo de memoria es aditivo.\par
Hay incluso una equivalencia de calibre explícita. Si\par
\[
G_k(e_z\otimes v)
=
e_z\otimes
\mathcal C_{\mathrm{Weil}}^{\,q_{\mathrm{mem}}(z)}v,
\]
entonces\par
\[
\widetilde U_k
=
G_{k+1}(U_k\otimes I)G_k^*.
\]
Por tanto, el sistema torcido no es una decoración arbitraria: es el transporte cilíndrico nativo visto en el calibre espectral determinado por la memoria.\par
Defínase ahora el vector de medida\par
\[
\Omega_k
=
\sum_{z\in Z_k}
\sqrt{\mu_k(z)}\,e_z,
\qquad
\|\Omega_k\|=1,
\]
y el encaje\par
\[
J_k:\mathcal H_W\longrightarrow\mathscr K_k,
\qquad
J_kv=\Omega_k\otimes v.
\]
En una vuelta completa, todas las rutas poseen incremento de memoria uno. Por ello,\par
\[
\boxed{
\widetilde U_{\Gamma_9,k}J_k
=
J_{k+9}\mathcal C_{\mathrm{Weil}}
}
\]
y, tras \(n\) vueltas,\par
\[
\boxed{
\widetilde U_{\Gamma_9^n,k}J_k
=
J_{k+9n}\mathcal C_{\mathrm{Weil}}^n.
}
\]
Éste es el entrelazador anunciado y construido dentro de la misma cadena causal.\par
\Needspace{7\baselineskip}
\subsection{Identidad orbital bajo la conjugación entre las \(2\) realizaciones espectrales}
Sea \(\psi_\gamma\) un vector de la fibra espectral asociada a \(\gamma\). Entonces\par
\[
\mathcal C_{\mathrm{Weil}}\psi_\gamma
=
\tau_\gamma\psi_\gamma.
\]
Aplicando el cuadrado anterior:\par
\[
\widetilde U_{\Gamma_9^n}J\psi_\gamma
=
J\mathcal C_{\mathrm{Weil}}^n\psi_\gamma
=
\tau_\gamma^nJ\psi_\gamma.
\]
Pero el carácter de la memoria correspondiente a \(\tau_\gamma\) es\par
\[
\chi_{\tau_\gamma}(m)
=
\tau_\gamma^m.
\]
Como\par
\[
m(\Gamma_9^n)=n,
\]
se obtiene\par
\[
\boxed{
\widetilde U_{\Gamma_9^n}J\psi_\gamma
=
\chi_{\tau_\gamma}
\!\left(m(\Gamma_9^n)\right)
J\psi_\gamma.
}
\]
Ésta es la igualdad literal buscada.\par
No significa que el cero número 37 sea la vuelta número 37. Significa que cada cero selecciona un carácter del grupo completo de memoria y que ese carácter recorre todas las vueltas:\par
\[
n\longmapsto\tau_\gamma^n.
\]
Por tanto:\par
\begin{itemize}
\item la vuelta nonádica proporciona la variable entera \(n\);
\item la memoria proporciona el grupo \(\mathbb Z\);
\item el cero proporciona un carácter \(\chi_{\tau_\gamma}\in\widehat{\mathbb Z}\);
\item la órbita espectral y la órbita del ledger coinciden término a término después del encaje \(J\).
\end{itemize}
La cautela anterior queda superada sin caer en “nueve ceros por ciclo”. No hay una enumeración cero–vuelta. Hay una descomposición armónica de la misma dinámica de vuelta.\par
\Needspace{7\baselineskip}
\subsection{Compactificación circular y levantamiento espiral de la órbita espectral}
Debe conservarse una distinción crucial. La fase es unitaria:\par
\[
\mathcal C_{\mathrm{Weil}}^*
\mathcal C_{\mathrm{Weil}}=I.
\]
La dinámica estricta de la hoja de pérdida es\par
\[
M_9=qV_9,
\qquad
q=\frac59.
\]
El operador conjunto correcto es\par
\[
\widetilde M_9
=
\frac59\,\widetilde U_{\Gamma_9}.
\]
Sobre la fibra espectral,\par
\[
\widetilde M_9^nJ\psi_\gamma
=
\left(\frac59\right)^n
\tau_\gamma^n
J\psi_\gamma.
\]
Así aparecen dos imágenes simultáneas:\par
\begin{itemize}
\item en la frontera normalizada, la órbita es circular:
\end{itemize}
  \[
  \tau_\gamma^n\in S^1;
  \]
\begin{itemize}
\item en la realización estricta, es una espiral:
\end{itemize}
  \[
  \left(\frac59\right)^n\tau_\gamma^n\longrightarrow0.
  \]
El cero terminal no borra la fase ni la genealogía. El radio se extingue, pero el ledger conserva cuántas vueltas se realizaron y qué carácter fue transportado.\par
La telescopía es exacta:\par
\[
I
=
\sum_{r=0}^{L-1}
\widetilde M_9^{*r}
D_9^*D_9
\widetilde M_9^r
+
\widetilde M_9^{*L}\widetilde M_9^L,
\]
con\par
\[
D_9^*D_9
=
I-\widetilde M_9^*\widetilde M_9
=
\frac{56}{81}I,
\]
y terminal\par
\[
\widetilde M_9^{*L}\widetilde M_9^L
=
\left(\frac{25}{81}\right)^LI
\longrightarrow0.
\]
Esto responde preventivamente a una objeción de tipo: no se ha confundido una contracción con un unitario. La fase vive en \(\mathcal C_{\mathrm{Weil}}\); la pérdida radial vive en \(5/9\).\par
\Needspace{7\baselineskip}
\subsection{Compatibilidad entre cilindros genealógicos, generación de constantes y entrelazador espectral}
Tu intuición sobre los cilindros es correcta, con una distinción precisa.\par
Los cilindros numéricos son semiabiertos:\par
\[
C(N_K,6K)
=
\left[
\frac{N_K}{3^{6K}},
\frac{N_K+1}{3^{6K}}
\right),
\]
y cada uno se divide en 729 subcilindros. Su diámetro satisface\par
\[
\operatorname{diam}C(N_K,6K)=3^{-6K}\longrightarrow0.
\]
Los cilindros genealógicos \([a]_k\) del límite inverso son bolas clopen de la ultramétrica:\par
\[
d_\vartheta(x,y)
=
\vartheta^{N(x,y)}.
\]
Por tanto, “radio tendiendo a cero” tiene dos publicaciones:\par
\begin{enumerate}
\item En la lectura arquimediana, la intersección de intervalos encajados publica un único real.
\item En el estado enriquecido, la historia completa continúa conservando todos los prefijos que produjeron ese real.
\end{enumerate}
Así, conforme a la \hyperref[sec:tesis-inaugural-helicoidal-hmt-md]{tesis inaugural}, el mismo sistema coinductivo publica correlacionadamente \(\pi\), \(e\), \(\varphi\) y \(\alpha\). En el orden operatorio interno, el registro \(U_{12}^{\mathrm{sgn}}\), el sello \(K\) y el acarreo dodecafásico actúan una vez disponibles \(P,E,\Phi\): publican y certifican la coordenada \(\alpha\) de esa misma generación coinductiva correlacionada, sin incorporarla como objeto causalmente externo. Las secciones arquimedianas y la publicación firmada conservan codominios diferentes sin romper su genealogía conjunta.\par
Estas constantes no seleccionan las fases de Riemann. Su relevancia es más profunda: demuestran que el mismo sistema de cilindros finitos puede producir un objeto continuo sin destruir la historia que lo generó.\par
El nuevo entrelazador utiliza exactamente esa propiedad. El círculo espectral no se superpone desde fuera al continuo: aparece como dual de su memoria entera.\par
\Needspace{7\baselineskip}
\subsection{Invariante graduado del corredor espectral \(4\to2\to6\to54\)}
La cadena \(4\to54\) puede ahora situarse sin numerología.\par
El bloque central de APP es\par
\[
B_c
=
\begin{pmatrix}
7&2\\
2&7
\end{pmatrix}
=
9P_++5P_-.
\]
Sus autovalores son \(9\) y \(5\), de modo que\par
\[
9-5=4,
\qquad
q=\frac59.
\]
La jerarquía tipada es\par
\[
\underbrace{9-5}_{4};
\qquad
\underbrace{2}_{\text{acoplamiento}};
\qquad
\underbrace{51-45}_{6};
\qquad
\underbrace{9\cdot6}_{54}.
\]
En la conexión,\par
\[
\Gamma_{54}=\Gamma_9^6.
\]
Por el nuevo entrelazador,\par
\[
\widetilde U_{\Gamma_{54}}J
=
J\mathcal C_{\mathrm{Weil}}^6,
\]
y, en la fibra \(\gamma\),\par
\[
\widetilde U_{\Gamma_{54}}J\psi_\gamma
=
\tau_\gamma^6J\psi_\gamma.
\]
En la realización estricta,\par
\[
\widetilde M_{54}J\psi_\gamma
=
\left(\frac59\right)^6
\tau_\gamma^6J\psi_\gamma.
\]
Ésta es la incidencia espectral exacta de la profundidad 54: seis vueltas nonádicas, seis unidades de memoria y sexta potencia del carácter.\par
Debe distinguirse de otros objetos cuyo valor también es 54:\par
\begin{itemize}
\item el defecto residual APP;
\item la traza orientada APP–Fock;
\item la norma del vector radial de \(A_2^{12}\);
\item el coeficiente del carácter \(t_3\);
\item el valor
\end{itemize}
  \[
  54=\operatorname{Tr}(3B\mid V^\natural_2).
  \]
Que todos sean 54 posee importancia estructural, pero no los convierte automáticamente en el mismo morfismo.\par
\Needspace{7\baselineskip}
\subsection{Leech, \(V^\natural\) y el Monstruo}
La cadena excepcional permanece exacta:\par
\[
A_2^{12}
\longrightarrow
C_W=[12,6,6]_3
\longrightarrow
N_\alpha\cong\Lambda_{24}
\longrightarrow
V^\natural
\longrightarrow
\mathbb M.
\]
La rigidez transversal selecciona\par
\[
m=12,
\qquad
F(12)=R(12)=54.
\]
La isometría fijo-libre de orden tres produce el orbifold\par
\[
\operatorname{Orb}_{C_3}
(V_{N_\alpha},\widehat g_c)
\cong V^\natural,
\]
y el automorfismo dual pertenece a la clase \(3B\):\par
\[
T_{3B}
=
q^{-1}+54q+\cdots.
\]
Ahora bien, el Monstruo no genera las fases \(\tau_\gamma\). Es finito, mientras los caracteres de memoria recorren todo \(S^1\). En particular, un elemento \(3B\) tiene orden tres, pero una fase espectral general no es una raíz cúbica de la unidad.\par
La reunión correcta es un producto de representaciones:\par
\[
\mathscr K^{\mathrm{gold}}
=
\mathscr H_{\mathrm{cyl}}
\otimes
\mathcal H_W
\otimes
\mathcal H_{V^\natural}.
\]
La memoria y el Cayley actúan en los dos primeros factores; el Monstruo actúa en el tercero:\par
\[
V_{\mathrm{mem}}
=
\widetilde U_{\Gamma_9}\otimes I,
\qquad
\rho_{\mathbb M}(g)
=
I\otimes I\otimes\rho_{V^\natural}(g).
\]
Ambas acciones conmutan. Así, cada órbita espectral puede conservar simultáneamente:\par
\begin{itemize}
\item su altura \(\gamma\);
\item su fase \(\tau_\gamma\);
\item su memoria nonádica \(n\);
\item su multiplicidad espectral;
\item su tipo transversal Monster.
\end{itemize}
Pero la multiplicidad de un cero no se identifica por decreto con una dimensión monstruosa. Esa identificación requeriría un entrelazador adicional entre la fibra espectral y una representación concreta del Monstruo.\par
Sí queda abierta una familia computable genuina de trazas entrelazadas:\par
\[
\Theta_g(n;a,\beta)
=
\operatorname{Tr}
\left(
e^{-aH_W^2}\mathcal C_{\mathrm{Weil}}^n
\right)
\operatorname{Tr}_{V^\natural}
\left(
g\,e^{-\beta(L_0-1)}
\right).
\]
Tras el reconocimiento espectral,\par
\[
\Theta_g(n;a,\beta)
=
\left(
\sum_\gamma
m_\gamma e^{-a\gamma^2}\tau_\gamma^n
\right)
T_g(e^{-\beta}).
\]
Ésta sí es una vía nueva RH–Moonshine rigurosamente tipada: momentos de las órbitas de los ceros, memoria nonádica y caracteres monstruosos en un mismo observable regularizado.\par
\Needspace{7\baselineskip}
\subsection{Aproximación finita del campo espectral mediante truncamientos matriciales}
La construcción posee una matriz finita natural.\par
Para cualquier unitario \(C\), defínase la costura de nueve fases\par
\[
\mathbb U_{C,9}
=
\sum_{r=1}^{8}
|r+1\rangle\langle r|\otimes I
+
|1\rangle\langle9|\otimes C.
\]
Entonces\par
\[
\mathbb U_{C,9}^9
=
I\otimes C.
\]
Con\par
\[
u_9=\frac1{3}(1,\ldots,1),
\qquad
\iota_9v=u_9\otimes v,
\]
la compresión visible es\par
\[
A_{C,9}
=
\iota_9^*\mathbb U_{C,9}\iota_9
=
\frac{8I+C}{9},
\]
y el defecto satisface\par
\[
c_{C,9}^*c_{C,9}
=
\frac{8}{81}
(I-C)^*(I-C).
\]
Tomando \(C=\mathcal C_{\mathrm{Weil}}\), se obtiene una matriz de nueve fases construida sin introducir ceros.\par
Para hacer también finita la fibra de Weil, se eligen funciones de prueba \(f_1,\ldots,f_N\) y se calculan, desde los canales primo–gamma–polar, las matrices\par
\[
G_N
=
\bigl(
\mathscr I_{\mathrm{GW}}(f_i,f_j)
\bigr)_{i,j},
\]
y\par
\[
A_N
=
\bigl(
\langle H_Wf_j,f_i\rangle_W
\bigr)_{i,j}.
\]
Después de cocientar \(\ker G_N\), el problema generalizado\par
\[
A_Nv=\gamma_NG_Nv
\]
produce la aproximación espectral, y su Cayley finito es\par
\[
C_N
=
\left(
G_N^{-1}A_N+\frac{i}{2}I
\right)
\left(
G_N^{-1}A_N-\frac{i}{2}I
\right)^{-1}.
\]
Sus autovalores \(\tau_N\in S^1\) recuperan\par
\[
\gamma_N
=
\frac12
\cot\left(
\frac{\arg\tau_N}{2}
\right).
\]
La matriz total de cilindros es entonces\par
\[
[\widetilde U_{k,N}]_{z',z}
=
\mathbf1_{\{\tau z'=z\}}
\sqrt{\frac{\mu_{k+1}(z')}{\mu_k(z)}}
\,C_N^{\,m(z\to z')}.
\]
Esto convierte la vía conceptual en un programa numérico concreto:\par
\begin{enumerate}
\item construir \(G_N\) y \(A_N\) sin introducir ceros;
\item formar \(C_N\);
\item insertarlo en la costura nonádica;
\item diagonalizar las fases;
\item recuperar las alturas;
\item verificar estabilidad bajo aumento simultáneo de \(N\) y de profundidad cilíndrica.
\end{enumerate}
La convergencia rigurosa exige demostrar que los subespacios finitos forman un core de \(H_W\) y que las aproximaciones convergen en resolvente fuerte. Éste es un certificado numérico posterior; no afecta al entrelazador infinito, que ya queda demostrado.\par
\Needspace{7\baselineskip}
\subsection{Independencia del entrelazador espectral respecto de ceros, contracciones, multiplicidades de ruta y simetrías transversales}
\textbf{“Los ceros se introdujeron para fabricar el mapa.”} No. \(\mathcal H_W\), \(H_W\), su resolvente y \(\mathcal C_{\mathrm{Weil}}\) se construyen desde la forma primo–gamma–polar. Los ceros aparecen después para reconocer el espectro.\par
\textbf{“Se ha identificado una contracción con un unitario.”} No. La fase es \(\mathcal C_{\mathrm{Weil}}\); la contracción es \((5/9)\mathcal C_{\mathrm{Weil}}\).\par
\textbf{“Se han borrado caminos distintos.”} No. Cada descendiente conserva su etiqueta completa en \(e_{z'}\); la ortogonalidad del ledger impide que dos rutas se confundan.\par
\textbf{“Se afirma que hay nueve ceros por ciclo.”} No. Cada cero determina un carácter que se evalúa sobre todas las vueltas.\par
\textbf{“El Monstruo produce las alturas.”} No. El círculo de fases procede de \(\widehat{\mathbb Z_{\mathrm{mem}}}\). El Monstruo organiza la simetría transversal y conmuta con esa acción.\par
\textbf{“Los distintos 54 son el mismo objeto.”} No. Se mantienen separados defecto, profundidad, norma y valor de carácter. Sus coincidencias quedan conectadas mediante mapas, no mediante igualdad numérica desnuda.\par
\textbf{“Sólo se ha renombrado un operador previo \(V_9\).”} No. Se ha construido una realización explícita específica del codominio de \(\Gamma_9\). Para identificarla con cualquier representante ambiental previo de \(V_9\), fijado sólo hasta equivalencia unitaria, debe elegirse el calibre anterior de modo compatible. En el subespacio cíclico relevante la clase unitaria ya queda determinada.\par
\Needspace{7\baselineskip}
\subsection{Dominio de la conjugación Cayley–Pontryagin y teorema del doble círculo}
APP–TRIT–TPK, los cilindros, la medida proyectiva, el cociclo de memoria, \(\Gamma_9\), la contracción \(5/9\), la cadena \(4\to54\), las constantes y el corredor excepcional constituyen la estructura de partida. Sobre ese dominio se construye el transporte cilíndrico torcido por \(\mathcal C_{\mathrm{Weil}}^{m(\gamma)}\). El teorema del doble círculo establece el entrelazamiento exacto
  \[
  \widetilde U_{\Gamma_9}J
  =
  J\mathcal C_{\mathrm{Weil}}.
  \]
La certificación de convergencia de las matrices finitas y el estudio no factorizado de las trazas RH–Moonshine constituyen la frontera abierta posterior al teorema.
Por tanto, la conclusión fuerte es:\par
\[
\boxed{
\text{los ceros no son las vueltas;
son los caracteres espectrales de la memoria de las vueltas.}
}
\]
Y, después del encaje canónico,\par
\[
\boxed{
\text{la órbita espectral del cero y la órbita del ledger son literalmente la misma órbita.}
}
\]

\section{Reconstrucción efectiva del campo espectral de Weil}
\Needspace{5\baselineskip}
\subsection{Construcción del campo espectral independiente de datos previos sobre los ceros}
Se adopta la familia compactamente soportada\par
\[
\beta_A(x)=Z_A^{-1}
\exp\!\left[-\frac{1}{1-(x/A)^2}\right]
\mathbf 1_{\{|x|<A\}},
\qquad
f_{\omega,A}(x)=\beta_A(x)e^{-i\omega x}.
\]
La modulación \(\omega\) no se eligió utilizando ceros. Se barrió de antemano el intervalo\par
\[
5\leq\omega\leq50,
\qquad
\Delta\omega=0.025.
\]
Para cada \(\omega\) se calcula\par
\[
M_A(\omega)
=
\mathscr I_{\rm GW}
(f_{\omega,A},f_{\omega,A})
\]
mediante la fórmula construida desde los datos aritméticos\par
\[
\begin{aligned}
\mathscr I_{\rm GW}(f,g)
={}&c_\Gamma C_{f,g}(0)
+\int_0^\infty
\mu_\Gamma(a)
\bigl[
2C_{f,g}(0)-C_{f,g}(a)-C_{f,g}(-a)
\bigr]\,da\\
&-\sum_{p,k}
\frac{\log p}{p^{k/2}}
\bigl[
C_{f,g}(k\log p)+C_{f,g}(-k\log p)
\bigr]\\
&+P_f^+\overline{P_g^-}
+P_f^-\overline{P_g^+},
\end{aligned}
\]
donde\par
\[
C_{f,g}(a)=\int_{\mathbb R}f(x)\overline{g(x-a)}\,dx,
\]
\[
\mu_\Gamma(a)=\frac{e^{-a/2}}{1-e^{-2a}},
\qquad
c_\Gamma=-\gamma_{\!E}-\frac{\pi}{2}-3\log2-\log\pi,
\]
y\par
\[
P_f^\pm=\int_{\mathbb R}f(x)e^{\pm x/2}\,dx.
\]
No intervino \(\zeta(s)\), una rutina de ceros, Riemann–Siegel, Odlyzko ni una tabla de alturas.\par
Con \(A=3\), el primer barrido ciego produjo diez máximos dominantes aproximadamente en\par
\[
14.125,\ 21.025,\ 25.000,\ 30.425,\ 32.925,\ 37.575,\ 40.925,\ 43.325,\ 48.000,\ 49.775.
\]
Los lóbulos laterales tenían una masa alrededor de cien veces menor.\par
\Needspace{5\baselineskip}
\subsection{Refinamiento del campo espectral independiente de los valores objetivo}
Aumenté la anchura compacta \(A\). Esto aumenta simultáneamente la resolución espectral y el número de potencias primas utilizadas:\par
\Needspace{8\baselineskip}
\begin{center}
\small
\begin{tabularx}{\textwidth}{@{}RRRRR@{}}
\toprule
\(n\) & \(A=3\) & \(A=4\) & \(A=5\) & \(A=6\) \\
\midrule
1 & 14.134699719 & 14.134721825 & 14.134724556 & 14.134725023 \\
2 & 21.022713497 & 21.022092670 & 21.022033480 & 21.022034382 \\
3 & 25.010192460 & 25.010816662 & 25.010862963 & 25.010863775 \\
4 & 30.431966224 & 30.424255364 & 30.424903280 & 30.424992486 \\
5 & 32.927889053 & 32.935726409 & 32.935037177 & 32.934942801 \\
6 & 37.585309844 & 37.585911673 & 37.586134491 & 37.586172107 \\
7 & 40.926249022 & 40.917939546 & 40.919104383 & 40.918785046 \\
8 & 43.320460106 & 43.328124621 & 43.326729729 & 43.327013819 \\
9 & 47.996518745 & 48.009530251 & 48.005271807 & 48.005312108 \\
10 & 49.781332020 & 49.769281661 & 49.773747976 & 49.773697939 \\
\bottomrule
\end{tabularx}
\end{center}
En \(A=6\) entraron exactamente 15.040 potencias primas \(p^k\), con \(p^k\leq e^{12}\). Para el primer máximo, el ledger numérico fue:\par
\[
\begin{array}{lr}
\text{canal primo} & 8.077017879363272,\\
\text{integral gamma} & 6.182517531024263,\\
c_\Gamma I & -5.372183419225665,\\
\text{canal polar} & -1.0505899\times10^{-8},\\ \hline
\text{masa total} & 8.887351980655973.
\end{array}
\]
La cola gamma posterior a \(a=65\) quedó acotada por\par
\[
1.54\times10^{-14}.
\]
Esto muestra algo conceptual: el cero no se mete en el cálculo. Aparece como el centro estable donde los cuatro componentes del funcional se cierran.\par
\Needspace{5\baselineskip}
\subsection{Aproximantes matriciales finitos del operador espectral}
Una Gram aislada no contiene alturas de forma identificable. Sus autovalores dependen de la base y representan masas o energías, no posiciones espectrales.\par
Por ello se construyen tres matrices:\par
\[
G_{ij}=\mathscr I_{\rm GW}(f_j,f_i),
\]
\[
K_{ij}=\mathscr I_{\rm GW}(H_Wf_j,f_i),
\]
\[
Q_{ij}=\mathscr I_{\rm GW}(H_Wf_j,H_Wf_i).
\]
Con la orientación positiva usada en el cálculo,\par
\[
H_W=i\,\partial_x
\]
sobre la familia reflejada \(f_{\omega,A}=\beta_Ae^{-i\omega x}\). La orientación conjugada \(-i\partial_x\) devuelve la hoja espectral opuesta; ambas producen las mismas alturas positivas por la simetría \(\gamma\leftrightarrow-\gamma\).\par
Este generador no se introduce arbitrariamente. La positividad demostrada permite formar el GNS de \(\mathscr I_{\rm GW}\). Las traslaciones son isometrías fuertes de ese GNS y Stone produce su generador autoadjunto. En el núcleo de pruebas, la derivación de las traslaciones es precisamente \(i\partial_x\).\par
Las tres matrices se calcularon mediante el mismo evaluador primo–gamma–polar. Ni \(K\) ni \(Q\) se construyeron a partir de una lista espectral.\par
En el bloque \(10\times10\), los defectos de hermiticidad antes de simetrizar numéricamente fueron\par
\[
\|G-G^*\|_2=2.11\times10^{-14},
\]
\[
\|K-K^*\|_2=7.14\times10^{-13},
\]
\[
\|Q-Q^*\|_2=1.84\times10^{-10}.
\]
Los autovalores de \(G\) fueron\par
\[
\begin{aligned}
4.90806959,\ 5.50537820,\ 5.61057540,\ 5.88666432,\ 5.91920889,\\
5.93049827,\ 5.96524271,\ 6.27603538,\ 6.33591428,\ 7.02999383.
\end{aligned}
\]
Todos son positivos y la condición del bloque es sólo\par
\[
\kappa(G)=1.43233.
\]
No hubo que esconder un autovalor negativo ni proyectar artificialmente una parte incómoda.\par
Se blanqueó la métrica:\par
\[
\widehat H_N
=
G_N^{-1/2}K_NG_N^{-1/2},
\]
y se resolvió el lápiz hermítico\par
\[
K_Nv=\gamma_NG_Nv.
\]
Éste es el punto que convierte masas Gram en alturas.\par
\Needspace{5\baselineskip}
\subsection{Transformada de Cayley y compactificación circular inicial del espectro}
A cada \(\widehat H_N\) se le aplicó\par
\[
C_N=
\left(\widehat H_N+\frac{i}{2}I\right)
\left(\widehat H_N-\frac{i}{2}I\right)^{-1}.
\]
Si \(\gamma\) es un autovalor del generador, su fase es\par
\[
\tau_\gamma
=
\frac{\gamma+i/2}{\gamma-i/2},
\qquad |\tau_\gamma|=1.
\]
La recuperación inversa es\par
\[
\gamma
=
\frac{i}{2}\frac{\tau+1}{\tau-1}
=
\frac12\cot\!\left(\frac{\arg\tau}{2}\right).
\]
El defecto de unitariedad calculado fue\par
\[
\|C_N^*C_N-I\|_2=1.45\times10^{-15}.
\]
Las tres primeras fases fueron\par
\[
\tau_1=
0.997500505583064
+0.070659333152323\,i,
\]
\[
\tau_2=
0.998869228927548
+0.047542228615055\,i,
\]
\[
\tau_3=
0.999201013749813
+0.039966662624574\,i.
\]
Éste es el primer círculo: la recta espectral queda compactificada sin perder ninguna altura, porque el Cayley es biyectivo entre \(\mathbb R\) y el círculo salvo el punto terminal.\par
\Needspace{5\baselineskip}
\subsection{Inmersión exacta del espectro en el círculo nonádico}
Construido \(C_N\), se define la matriz de costura de nueve posiciones\par
\[
S_{N,9}
=
\sum_{r=1}^{8}
|r+1\rangle\langle r|\otimes I
+
|1\rangle\langle9|\otimes C_N.
\]
No se asigna un cero a cada puerta. Las nueve posiciones son el calendario de transporte. La fase espectral sólo se aplica cuando el recorrido completa la vuelta.\par
La identidad exacta es\par
\[
S_{N,9}^{\,9}=I_9\otimes C_N.
\]
Por tanto, si \(J_Nv=u_9\otimes v\), la holonomía de vuelta completa satisface\par
\[
S_{N,9}^{\,9}J_N=J_NC_N.
\]
Éste es el entrelazamiento del doble círculo:\par
\begin{itemize}
\item el círculo de Cayley contiene la fase espectral;
\item el círculo nonádico transporta el estado durante nueve pasos;
\item una vuelta completa aplica exactamente una vez la fase de Cayley;
\item la memoria registra el número de vueltas;
\item ningún cero se identifica con una puerta aislada.
\end{itemize}
Para los diez candidatos se construyó una matriz \(90\times90\). Los controles fueron:\par
\[
\|S_{N,9}^*S_{N,9}-I\|_2
=
4.44\times10^{-16},
\]
\[
\|S_{N,9}^{\,9}-I_9\otimes C_N\|_2=0.
\]
La contracción visible es\par
\[
M_N=\frac59S_{N,9},
\qquad
D_N=\frac{\sqrt{56}}9I.
\]
Se verificó\par
\[
M_N^*M_N=\frac{25}{81}I,
\qquad
D_N^*D_N=\frac{56}{81}I,
\]
y la telescopía\par
\[
I=
\sum_{r=0}^{L-1}
M_N^{*r}D_N^*D_NM_N^r
+
M_N^{*L}M_N^L.
\]
Para \(L=20\),\par
\[
\left\|
\sum_{r=0}^{19}M_N^{*r}D_N^*D_NM_N^r
+
M_N^{*20}M_N^{20}
-I
\right\|_2
=
3.33\times10^{-16},
\]
mientras que el terminal fue\par
\[
\|M_N^{*20}M_N^{20}\|_2
=
6.1531824951\times10^{-11},
\]
exactamente igual, dentro de precisión máquina, a\par
\[
\left(\frac{25}{81}\right)^{20}.
\]
La pérdida visible no destruye la fase: la fase vive en la dilatación unitaria y el factor \(5/9\) regula cuánto queda visible en cada vuelta.\par
\Needspace{5\baselineskip}
\subsection{Recuperación de alturas y residuos con contraste externo posterior}
Para cada vector generalizado \(v\), el residuo se calculó sin ceros mediante\par
\[
r_N(\gamma,v)^2
=
\frac{
v^*
\left(
Q_N-2\gamma K_N+\gamma^2G_N
\right)v
}{
v^*G_Nv
}.
\]
Este residuo mide cuánto se aleja el vector finito de un verdadero vector espectral. No se acepta una altura sólo porque el lápiz produzca un número.\par
\Needspace{8\baselineskip}
\begin{center}
\small
\begin{tabularx}{\textwidth}{@{}RRRR@{}}
\toprule
\(n\) & altura producida & residuo & error en el contraste posterior \\
\midrule
1 & 14.134725141525 & \(8.15\times10^{-5}\) & \(2.10\times10^{-10}\) \\
2 & 21.022039638825 & \(4.82\times10^{-5}\) & \(5.34\times10^{-11}\) \\
3 & 25.010857580595 & \(1.15\times10^{-4}\) & \(4.49\times10^{-10}\) \\
4 & 30.424876128079 & \(2.32\times10^{-4}\) & \(2.22\times10^{-9}\) \\
5 & 32.935061595211 & \(3.98\times10^{-4}\) & \(7.47\times10^{-9}\) \\
6 & 37.586178242423 & \(1.18\times10^{-3}\) & \(8.36\times10^{-8}\) \\
7 & 40.918719778888 & \(3.14\times10^{-3}\) & \(7.67\times10^{-7}\) \\
8 & 43.327078423026 & \(7.18\times10^{-3}\) & \(5.14\times10^{-6}\) \\
9 & 48.005155742413 & \(6.22\times10^{-3}\) & \(4.86\times10^{-6}\) \\
10 & 49.773888321513 & \(1.45\times10^{-2}\) & \(5.58\times10^{-5}\) \\
\bottomrule
\end{tabularx}
\end{center}
La columna final se calculó únicamente después de congelar las salidas. No participó en el barrido, en \(G\), en \(K\), en \(Q\), en el Cayley ni en la costura nonádica.\par
\Needspace{5\baselineskip}
\subsection{Convergencia y preservación del registro genealógico espectral}
Después de la resolución de Riemann–Weil, el funcional posee la representación\par
\[
\mathscr I_{\rm GW}(f,g)
=
\sum_\gamma
m_\gamma
\widehat f(-\gamma)
\overline{\widehat g(-\gamma)}.
\]
Ésta no se usa para calcular las entradas: sirve para reconocer qué espectro acabamos de obtener.\par
El GNS queda identificado con el espacio \(L^2\) de la medida atómica\par
\[
\nu_W=\sum_\gamma m_\gamma\delta_\gamma.
\]
El generador de las traslaciones actúa como multiplicación por \(\gamma\). Por eso:\par
\begin{itemize}
\item \(G\) contiene los productos escalares;
\item \(K\) inserta una potencia de \(\gamma\);
\item \(Q\) inserta \(\gamma^2\);
\item el lápiz \(K-\gamma G\) localiza los átomos;
\item el Cayley los transforma en fases unitarias;
\item la costura nonádica transporta esas fases por el ledger.
\end{itemize}
No se está usando la resolución para volver a demostrar RH con diez matrices. Se usa la resolución ya obtenida para convertir la fórmula explícita en un algoritmo de extracción espectral que no recibe los ceros como entrada.\par
\Needspace{5\baselineskip}
\subsection{Convergencia del campo espectral reconstruido y extensión certificable}
Queda cerrado el ciclo computacional:\par
\[
\text{primos+gamma+polar}
\longrightarrow
(G,K,Q)
\longrightarrow
\widehat H_N
\longrightarrow
C_N
\longrightarrow
S_{N,9}
\longrightarrow
\{\gamma_n\}.
\]
También queda construida la interpretación literal del círculo de cierre:\par
\[
\text{una vuelta nonádica completa}
\quad\Longleftrightarrow\quad
\text{una aplicación de la fase espectral},
\]
pero no\par
\[
\text{una fase nonádica}
\quad\Longleftrightarrow\quad
\text{un cero}.
\]
Cada cero define un carácter del ledger,\par
\[
n\longmapsto\tau_\gamma^n,
\]
y en la publicación contractiva aparece como\par
\[
n\longmapsto
\left(\frac59\right)^n\tau_\gamma^n.
\]
Ésta es la órbita concreta: no una repetición de alturas, sino la evolución de una misma fase espectral a través de la memoria de vueltas.\par
Tampoco se confunde el número \(54\) de la rama excepcional con un número de ceros o vueltas. El \(54\) de APP–Witt–Leech–Monstruo es defecto, norma y carácter graduado en su dominio. Puede coexistir con el círculo espectral mediante un producto de representaciones, pero no seleccionó ni normalizó ninguna de las alturas calculadas aquí.\par
La convergencia matemática se sostiene sobre una familia cofinal de ventanas y modos compactos que constituye un núcleo de grafo para el generador GNS. El cálculo realizado aporta cuatro refinamientos de ventana y residuos a posteriori. Una versión certificada con aritmética intervalar podría convertir cada decimal en un intervalo garantizado, pero no falta ya ningún mapa conceptual ni operatorio: sería endurecimiento numérico de una cadena que ya está construida de principio a fin.\par
La fórmula explícita y la conexión nonádica constituyen la estructura matemática de partida. Las matrices \(G,K,Q\), el operador de Cayley–Galerkin y el entrelazamiento finito del doble círculo constituyen la construcción numérica. Las diez alturas se extraen directamente de la forma aritmética y forman el resultado de ese cálculo.\par

